% =========================================================
%  SECCIÓN I — INTRODUCCIÓN
% =========================================================
\section{Introducci\'on}

\PARstart{L}{a} detecci\'on de bordes es una de las operaciones
fundamentales en el procesamiento digital de im\'agenes.
Su objetivo consiste en identificar los cambios abruptos en la
intensidad luminosa de los p\'ixeles de una imagen, los cuales
generalmente corresponden a las fronteras entre objetos, texturas
o regiones de inter\'es~\cite{gonzalez2008}.
Esta operaci\'on es el primer paso en cadenas de an\'alisis m\'as
complejas, como la segmentaci\'on sem\'antica, el reconocimiento
de objetos o la reconstrucci\'on tridimensional en aplicaciones
de visi\'on artificial, medicina, rob\'otica y teledetección
satelital~\cite{canny1986}.

El rendimiento de estas operaciones se vuelve cr\'itico cuando
se procesan grandes vol\'umenes de informaci\'on, como video en
tiempo real (30--60 cuadros por segundo), im\'agenes m\'edicas de
alta resoluci\'on (tomograf\'ia, resonancia magn\'etica) o flujos
de datos provenientes de sensores en sistemas embebidos.
Los procesadores modernos incorporan unidades de ejecuci\'on
vectorial (SIMD, \textit{Single Instruction, Multiple Data})
capaces de procesar m\'ultiples datos en paralelo dentro de un
\'unico n\'ucleo de c\'omputo.
En la arquitectura x86-64, la extensi\'on AVX2 (\textit{Advanced
Vector Extensions 2}) permite operar con registros YMM de 256~bits,
procesando hasta ocho valores de punto flotante simple
(32~bits) en una sola instrucci\'on~\cite{intel_intrinsics}.

Sin embargo, explotar correctamente la vectorizaci\'on SIMD no es
trivial.
El compilador puede vectorizar autom\'aticamente cuando detecta
patrones seguros, pero en ocasiones requiere ayuda mediante pragmas
o directivas.
Para m\'axima eficiencia, el programador puede emplear intr\'insecos
AVX2 directamente, aunque esto eleva la complejidad del c\'odigo.
Independientemente de la estrategia de vectorizaci\'on adoptada,
el rendimiento final puede estar limitado por el ancho de banda
de memoria del sistema, un fen\'omeno conocido como
\textit{memory-bound}~\cite{williams2009}.

El presente trabajo analiza estas cuatro estrategias aplicadas
al kernel num\'erico del operador Prewitt~\cite{gonzalez2008,prewitt1970}.
Se elige el operador Prewitt frente al Sobel por su estructura
con pesos uniformes ($\pm 1$), que produce una intensidad
aritm\'etica te\'orica reducida ($\mathcal{I} \approx 0.333$~FLOP/Byte)
y lo convierte en un candidato ideal para estudiar el
comportamiento \textit{memory-bound} en el modelo Roofline.
Se utiliza Intel Advisor como herramienta de perfilado para
obtener datos objetivos sobre intensidad aritm\'etica, modelo
Roofline y comportamiento de los bucles vectorizados.
La hip\'otesis de trabajo es que la intensidad aritm\'etica del
kernel ($\mathcal{I} \approx 0.333$ FLOP/Byte) situar\'a todas
las versiones en la zona de memoria limitada del modelo Roofline,
independientemente del nivel de vectorizaci\'on aplicado.

El resto del art\'iculo est\'a organizado de la siguiente manera:
la Secci\'on~II presenta el marco te\'orico sobre el operador
Prewitt, vectorizaci\'on SIMD e Intel Advisor;
la Secci\'on~III describe la metodolog\'ia experimental;
la Secci\'on~IV presenta los resultados obtenidos;
la Secci\'on~V discute e interpreta los datos;
y la Secci\'on~VI expone las conclusiones del trabajo.
